%%%% RESUMO
%%
%% Apresentação concisa dos pontos relevantes de um texto, fornecendo uma visão rápida e clara do conteúdo e das conclusões do
%% trabalho.

\begin{resumoutfpr}%% Ambiente resumoutfpr
    Este estudo investiga os efeitos do \textit{fine-tuning} em modelos de linguagem encarregados de navegar por labirintos em ASCII por meio da utilização da biblioteca LLM-MRI. Ao analisar os grafos de ativações gerados a partir dos embeddings dos modelos, avaliamos as mudanças estruturais utilizando diversas métricas de grafos com inferência de distribuições por amostragem estatística. Nossos achados indicam alterações sutis, porém não significativas, nas estruturas dos grafos após o \textit{fine-tuning}, sugerindo que, embora o \textit{fine-tuning} influencie o processamento de informações, essas mudanças não são drasticamente refletidas nas métricas analisadas. O estudo destaca a importância de considerar diferentes níveis de abstração dentro dos modelos para compreender plenamente o impacto do \textit{fine-tuning}. Trabalhos futuros recomendados incluem a exploração de prompts diversos, uso de \textit{embbeddings} que capturem a resolução da tarefa e os efeitos de várias técnicas de \textit{fine-tuning} e arquiteturas de modelos nas estruturas dos grafos de ativações.
\end{resumoutfpr}
